\chapter*{Abstract}
\addcontentsline{toc}{chapter}{Abstract}

Predicting the biological function of proteins at scale remains one of the central
challenges in computational biology. Experimental characterisation cannot keep pace
with the rate at which new sequences are deposited in public databases, leaving the
vast majority of known proteins with no experimentally validated functional annotation.
The gap is widening: large-scale protein language models (PLMs) and structural
predictors generate dense representations of every UniProt sequence in days, yet the
infrastructure required to turn those representations into auditable, reproducible Gene
Ontology predictions is fragmented across notebooks, prototype scripts and one-off
benchmarks.

This thesis presents \textbf{PROTEA} (\textit{PROtein funcTional annotation framEwork
and plAtform}), an open-source distributed platform that automates the full pipeline
from raw protein sequence to structured Gene Ontology predictions, and provides an
extensible foundation on top of which third parties can plug in new PLMs, annotation
sources, and ranking methods. PROTEA combines protein language model embeddings (eight
backends evaluated, including ESM-2 in three sizes, ESM-C, ProstT5, ProtT5, Ankh and
ESM-3 component encoder), approximate $k$-nearest-neighbour search over 527,000 reviewed
UniProt sequences, and evidence transfer from versioned annotation sets to produce
ranked, interpretable predictions. A versioned relational data model links every
prediction to the exact ontology release, annotation set, embedding configuration and
code commit that produced it, enabling exact replay of any past run.

To improve prediction quality beyond embedding cosine similarity, PROTEA incorporates a
plugin-based feature engineering layer that computes Needleman--Wunsch and
Smith--Waterman alignment statistics, NCBI-taxonomy-based distance metrics, ancestry
embeddings of the GO directed acyclic graph, and per-PLM PCA projections, yielding a
schema of approximately 52 features per candidate annotation. A LightGBM re-ranker is
trained on these features using a temporal holdout protocol spanning twelve consecutive
GOA releases, with Information Accretion weighting, per-category models (No-Knowledge,
Limited-Knowledge, Prior-Knowledge), and selective application that falls back to the
embedding-only baseline when re-ranking would degrade performance.

Beyond the prediction method, the thesis describes the engineering pillars that make
the platform deployable across heterogeneous environments: a seven-repository plugin
architecture discovered via Python entry points, externalised configuration via
\texttt{pydantic-settings}, OpenTelemetry-based distributed tracing, multi-target
container packaging covering development, cloud, HPC (Apptainer) and airgapped
environments, and an operational narrative model that makes every Job and
\texttt{ExperimentRun} self-describing.

Evaluation follows the CAFA temporal protocol on held-out GOA splits and reports the
performance of the canonical pipeline across the eight PLM backends and three category
slices, contextualises the impact of each feature family through ablation studies, and
documents the methodological pivots that shaped the canonical configuration. Three
pre-built containers are submitted to the LAFA benchmark to demonstrate adoption-ready
deployment.

\bigskip
\noindent\textbf{Keywords:} protein function prediction, gene ontology, protein
language models, sequence embeddings, nearest-neighbour search, learning-to-rank,
LightGBM, distributed systems, reproducibility, container deployment, bioinformatics
infrastructure.
